% Kamp 1968, Chapter IV, Theorem 1 and its three lemmas, printed pp. 95--99.

\subsection{Why unary real-order tenses cannot express SINCE or UNTIL}
\label{sec:p09-kamp-iv1-unary-separation}

The expressive-completeness theorem permits binary temporal operations.
Kamp's next result shows that their arity cannot be reduced to one, even
when every first-order-definable unary tense on the real order is available.
The proof has three steps: one discrete proposition generates only four
truth sets; finitely many disjoint such propositions generate a finite
Boolean algebra; every formula using unary tenses preserves that algebra.
An even/odd-integer interpretation then separates this language from each
of the strict binary operations \(S,U\)
\cite[ch.~IV, thm.~1 and lemmas~1--3, pp.~95--99]{Kamp1968}.

\begin{hbdefinition}{def:p09-kamp-real-unary-language}{The quantified class of unary real-order tenses}
Throughout this section the time structure is exactly
\(\mathcal R=\langle\mathbb R,<\rangle\). A proposition is a subset
of \(\mathbb R\). For a subset \(A\), write
\(A^c=\mathbb R\setminus A\); complements are always relative to
this fixed time domain.

Let \(\mathfrak M_{\mathcal R}\) be the set of all operations
\(O:\mathcal P(\mathbb R)\to\mathcal P(\mathbb R)\) defined by
a source \(L\)-formula \(\varphi(t_0)\) whose only free variable
is \(t_0\), which contains \(P_0\), and which contains no
other unary predicate symbol:
\[
 O(A)=\{t\in\mathbb R:
       \langle\mathbb R,<,A\rangle\models\varphi[t_0\mapsto t]\}.
\]
The signature contains no real constants or additional predicate
parameters. The required occurrence of \(P_0\) may be
inessential: conjoining \(P_0(t_0)\lor\neg P_0(t_0)\)
inserts one without changing the defined operation. The operation
need only have the displayed definition on \(\mathcal R\).

Write \(\mathsf{TM}_{\mathcal R}\) for the language of finite
formulas generated by
\[
 \eta ::= q_i\mid\neg\eta\mid(\eta\land\eta)
                 \mid O(\eta),\qquad O\in\mathfrak M_{\mathcal R}.
\]
For a valuation \(V\) assigning propositions to letters, let
\(\llbracket\eta\rrbracket_V\) denote the truth set of \(\eta\).
In particular,
\(\llbracket O(\eta)\rrbracket_V
 =O(\llbracket\eta\rrbracket_V)\).
The alphabet contains all the permitted unary operations at once;
each individual formula is finite. Arbitrary finite truth-functional
sentential connectives may also be used, since negation and conjunction
express them. The constants are expressible using
\(q_0\land\neg q_0\) and its negation.

For propositions \(A,B\), the binary target operations have the
source's witness-first, open-interval semantics:
\[
\begin{aligned}
 S(A,B)&=\{t:\exists s<t\ (s\in A\land(s,t)\subseteq B)\},\\
 U(A,B)&=\{t:\exists s>t\ (s\in A\land(t,s)\subseteq B)\}.
\end{aligned}
\]
These are the same strict operations used in
\cref{thm:p09-kamp-iii2}. The quantification over unary tenses below
is over \(\mathfrak M_{\mathcal R}\), with the stated first-order
definitions and without additional parameters
\cite[ch.~II, defs.~2--3, pp.~21--22; ch.~IV, p.~95]{Kamp1968}.
\end{hbdefinition}

\begin{hbdefinition}{def:p09-kamp-du}{Discrete and unbounded real propositions}
A proposition \(A\subseteq\mathbb R\) is \emph{discrete} in
Kamp's sense when
\[
 \forall t\in\mathbb R\ \exists z,y\in\mathbb R\quad
 z<t<y,\qquad (z,t)\cap A=(t,y)\cap A=\varnothing.
\]
The condition is imposed at every real point, including points outside
\(A\). It excludes finite accumulation points.

The source calls \(A\) \emph{unbounded} when either
\(A=\varnothing\), or
\[
 \forall t\in\mathbb R\ \exists z,y\in A\quad z<t<y.
\]
Thus every nonempty unbounded proposition is unbounded in both
directions, and the empty proposition is expressly included.
Write \(\operatorname{DU}(A)\) when either \(A\) is
discrete and unbounded or its complement \(A^c\) is discrete
and unbounded, and put
\[
                  \mathcal F(A)=\{\varnothing,\mathbb R,A,A^c\}.
\]
The family \(\mathcal F(A)\) can have two members when
\(A=\varnothing\) or \(A=\mathbb R\)
\cite[ch.~IV, def.~1, p.~96]{Kamp1968}.
\end{hbdefinition}

\begin{hblemma}{thm:p09-kamp-iv1-lemma1}{proved=annex:p9.kamp.iv1.lemma1}{The four-set lemma for a single proposition}
Let \(\operatorname{DU}(A)\). If \(\eta\) is a
\(\mathsf{TL}_1\)-formula and every proposition letter occurring
in \(\eta\) is interpreted by \(A\), then
\[
                    \llbracket\eta\rrbracket\in\mathcal F(A).
\]
More generally the same conclusion holds when every occurring letter
is interpreted by any member of \(\mathcal F(A)\). The latter
form records the closure used in the induction; the former is Kamp's
Lemma~1. Both \(S\) and \(U\) send any two members of
\(\mathcal F(A)\) to either \(\varnothing\) or
\(\mathbb R\)
\cite[ch.~IV, lemma~1, pp.~96--97]{Kamp1968}.
\end{hblemma}

\begin{hbcorollary}{thm:p09-kamp-iv1-unary-corollary}{proved=annex:p9.kamp.iv1.unary.corollary}{Every definable unary tense preserves the four choices}
For every \(O\in\mathfrak M_{\mathcal R}\) and every
\(A\subseteq\mathbb R\),
\[
      \operatorname{DU}(A)\quad\Longrightarrow\quad
      O(A)\in\{\varnothing,\mathbb R,A,A^c\}.
\]
The choice among the four sets may depend on \(O\) and on \(A\).
This is a statement about every permitted unary tense, irrespective
of the number or alternation of quantifiers in its defining formula
\cite[ch.~IV, corollary to lemma~1, p.~97]{Kamp1968}.
\end{hbcorollary}

\begin{hbdefinition}{def:p09-kamp-discrete-generated-field}{The field generated by finitely many propositions}
For \(n\geq0\) and \(A_1,\ldots,A_n\subseteq\mathbb R\),
let \(\mathcal B(A_1,\ldots,A_n)\) be the smallest family of
subsets of \(\mathbb R\) containing these sets,
\(\varnothing\), and \(\mathbb R\), and closed under
complement and intersection. It is a Boolean algebra of sets, called
a field of subsets in the source. For
\(I\subseteq\{1,\ldots,n\}\), put
\[
                         A_I=\bigcup_{i\in I}A_i,
                  \qquad A_{\varnothing}=\varnothing.
\]
When \(n=0\), the generated field is
\(\{\varnothing,\mathbb R\}\).
\end{hbdefinition}

\begin{hblemma}{thm:p09-kamp-iv1-lemma2}{proved=annex:p9.kamp.iv1.lemma2}{Normal form in the finite generated field}
Let \(n\geq0\), and let \(A_1,\ldots,A_n\) be pairwise
disjoint, discrete, unbounded real propositions. Then
\[
\begin{aligned}
 \mathcal B(A_1,\ldots,A_n)
 &=\{A_I:I\subseteq\{1,\ldots,n\}\}\\
 &\quad\cup\{A_I^c:I\subseteq\{1,\ldots,n\}\}.
\end{aligned}
\]
In particular every member of the generated field satisfies
\(\operatorname{DU}\). Empty generators and an empty family are
allowed. The normal form is not asserted to be unique
\cite[ch.~IV, lemma~2, pp.~97--98]{Kamp1968}.
\end{hblemma}

\begin{hblemma}{thm:p09-kamp-iv1-lemma3}{proved=annex:p9.kamp.iv1.lemma3}{Closure under every unary-tense formula}
Under the hypotheses of \cref{thm:p09-kamp-iv1-lemma2}, let
\(\eta\in\mathsf{TM}_{\mathcal R}\) use only the letters
\(q_1,\ldots,q_n\), interpreted respectively by
\(A_1,\ldots,A_n\). Then
\[
              \llbracket\eta\rrbracket_V
                    \in\mathcal B(A_1,\ldots,A_n).
\]
The same conclusion holds if further letters occur and each of their
interpretations already belongs to this generated field. Consequently,
if two real points have the same membership in every generator, every
such formula has the same truth value at those two points.
The condition concerns every letter actually occurring in the formula
\cite[ch.~IV, lemma~3, p.~98]{Kamp1968}.
\end{hblemma}

\begin{hbtheorem}{thm:p09-kamp-iv1}{proved=annex:p9.kamp.iv1}{Kamp's real-order unary separation theorem}
Neither SINCE nor UNTIL is expressible on
\(\langle\mathbb R,<\rangle\) using sentential connectives and
arbitrary unary real-order tenses. Precisely, for each
\(T\in\{S,U\}\), there is no finite formula
\(\eta\in\mathsf{TM}_{\mathcal R}\), with no proposition
letters other than \(q_0,q_1\), such that
\[
 \forall A,B\subseteq\mathbb R\quad
       \llbracket\eta\rrbracket_{q_0\mapsto A,\,q_1\mapsto B}
                           =T(A,B).
\]
The impossibility holds with the full alphabet
\(\mathfrak M_{\mathcal R}\) available, and therefore with any
chosen family of its operations, finite or infinite.

For the fixed interpretation
\[
 E=\{2k:k\in\mathbb Z\},\qquad
 H=\{2k+1:k\in\mathbb Z\},\qquad
 q_0\mapsto E,\quad q_1\mapsto H,
\]
all \(\mathsf{TM}_{\mathcal R}\)-formulas in these two letters
agree at every pair of noninteger points. In contrast,
\[
\begin{array}{c|cc}
 & -\tfrac12 & +\tfrac12\\ \hline
 S(q_0,\neg q_1)&\text{false}&\text{true}\\
 U(q_0,\neg q_1)&\text{true}&\text{false}
\end{array}
\]
This supplies the two distinct nonexpressibility conclusions
\cite[ch.~IV, thm.~1, pp.~95--99]{Kamp1968}.
\end{hbtheorem}

\begin{hbexample}{def:p09-kamp-iv1-integer-boundaries}{The complete truth sets in the separating interpretation}
For the \(E,H\) in \cref{thm:p09-kamp-iv1}, the generated
field consists of exactly
\[
 \varnothing,\quad\mathbb R,\quad E,\quad H,\quad\mathbb Z,
 \quad E^c,\quad H^c,\quad\mathbb R\setminus\mathbb Z.
\]
Its three atoms are \(E\), \(H\), and
\(\mathbb R\setminus\mathbb Z\). The binary truth sets are
\[
 S(E,H^c)=\bigcup_{k\in\mathbb Z}(2k,2k+1],\qquad
 U(E,H^c)=\bigcup_{k\in\mathbb Z}[2k-1,2k).
\]
Each contains some noninteger points and excludes others. At an even
integer both formulas are false; at an odd integer both are true.
The closed endpoint at an odd integer is required because the current
point lies outside the open interval on which the second operand is
tested. The annex proves these full descriptions along with the theorem.
\end{hbexample}

\begin{hbexample}{def:p09-kamp-iv1-du-boundaries}{The empty case and the role of both unbounded directions}
Both \(\varnothing\) and \(\mathbb R\) satisfy
\(\operatorname{DU}\), and their four-choice family reduces to
\(\{\varnothing,\mathbb R\}\). Thus every permitted unary
tense sends either of these two inputs to one of these two outputs.

The singleton \(A=\{0\}\) is discrete but fails the source's
unboundedness condition. The permitted unary strict-past tense
\[
 O(A)=\{t:\exists s<t\ (s\in A)\}
\]
has value \((0,\infty)\), which is none of
\(\varnothing,\mathbb R,\{0\},\mathbb R\setminus\{0\}\).
This demonstrates why the four-choice assertion cannot be extended to
all discrete propositions. Likewise, requiring isolation only at members
would be insufficient to state the source definition: the set
\(\{1/m:m\in\mathbb Z,\ m\geq1\}\) has isolated members
but violates its right-hand gap condition at zero.
\end{hbexample}
